\documentclass[11pt]{article}
\usepackage{amsmath,amssymb,amsthm}
\usepackage[margin=1.1in]{geometry}
\usepackage{booktabs}
\usepackage{hyperref}

\newtheorem{definition}{Definition}
\newtheorem{proposition}{Proposition}
\newtheorem{theorem}{Theorem}
\newtheorem{remark}{Remark}

\title{Markov Chains:\\
The Probability Behind $n$-grams, and Maximum Likelihood on
Dependent Data}
\author{Yuval Lavie}
\date{\today}

\begin{document}
\maketitle

\begin{abstract}
The LLM track's models all stand on one probabilistic object this
library never formalized: the Markov chain. This page defines it,
proves the reduction that makes $k$-th order chains (equivalently,
$n$-gram models) first-order objects, states Chapman--Kolmogorov,
stationarity, mixing, and the ergodic theorem --- and then does the
statistics: maximum likelihood for a chain observed along a path.
This is the library's first MLE on \emph{dependent} data, and the
bernoulli estimator report card survives with new proofs: the ergodic
theorem replaces the law of large numbers, row sample sizes become
random, and the variance formula gains an occupancy weight. Two
companion scripts verify every claim:
\texttt{markov\_chains\_demo.py} (the probability) and
\texttt{markov\_mle.py} (the statistics, on a weather chain with
nothing linguistic about it).
\end{abstract}

\section{The object}

\begin{definition}[Markov chain]
A sequence $X_1, X_2, \dots$ on a finite state space
$S = \{1,\dots,V\}$ is a \emph{homogeneous Markov chain} if the future
depends on the past only through the present:
\begin{equation}
  \Pr(X_{t+1} = j \mid X_t = i,\, X_{t-1}, \dots, X_1)
  \;=\; \Pr(X_{t+1} = j \mid X_t = i) \;=\; P_{ij},
  \label{eq:markovprop}
\end{equation}
independent of $t$. The $V \times V$ matrix $P$ is row-stochastic:
each row is a conditional distribution --- \emph{one multinoulli per
current state}. With an initial distribution $\pi_0$, the path
probability factorizes:
\begin{equation}
  \Pr(x_1, \dots, x_T)
  \;=\; \pi_0(x_1) \prod_{t=1}^{T-1} P_{x_t,\, x_{t+1}}.
  \label{eq:path}
\end{equation}
\end{definition}

\begin{definition}[$k$-th order chain]
$X_{t+1}$ may depend on the last $k$ states:
$\Pr(X_{t+1} = j \mid X_t, \dots, X_1) =
p(j \mid X_{t-k+1}, \dots, X_t)$. This is exactly the $n$-gram model
of the LLM track with $n = k + 1$.
\end{definition}

\begin{proposition}[Reduction to first order]
A $k$-th order chain on $S$ is a first-order chain on the state space
$S^k$ of $k$-tuples: set $Y_t = (X_{t-k+1}, \dots, X_t)$; then $Y$ is
Markov with transitions
\begin{equation}
  \Pr\bigl(Y_{t+1} = (x_{t-k+2}, \dots, x_t, j) \mid Y_t\bigr)
  \;=\; p(j \mid Y_t),
  \label{eq:lift}
\end{equation}
and zero probability to any tuple that does not shift correctly.
\end{proposition}

\begin{proof}
$Y_{t+1}$ is a deterministic shift of $Y_t$ plus the single new symbol
$X_{t+1}$, whose law depends only on the last $k$ states --- i.e.\
only on $Y_t$. Hence $Y$ satisfies \eqref{eq:markovprop} on $S^k$.
\end{proof}

This proposition is why ``$n$-gram model'' and ``Markov chain'' are the
same mathematics, and why the $k$-gram count table has $V^k$ rows (the
curse measured in \texttt{Generative/text-to-text/ngram-mlp}). Verified numerically on
real data: the trigram ($k=2$) log-likelihood of the Shakespeare
corpus computed directly and via the lifted pair-state chain agree to
$10^{-6}$ absolute --- term for term the same sum.

\section{Dynamics: $n$ steps, stationarity, mixing, ergodicity}

\begin{proposition}[Chapman--Kolmogorov]
$\Pr(X_{t+n} = j \mid X_t = i) = (P^n)_{ij}$: condition on the
intermediate state and sum --- which is exactly matrix multiplication.
Measured on a simulated path ($T = 2\times10^5$): empirical two-step
frequencies match $P^2$ to $0.004$.
\end{proposition}

\begin{definition}[Stationary distribution]
$\pi$ with $\pi P = \pi$ --- a left eigenvector of eigenvalue $1$,
normalized to a distribution. If the chain starts at $\pi$, every
marginal stays $\pi$.
\end{definition}

\begin{theorem}[Convergence and mixing]
If $P$ is irreducible (every state reachable from every state) and
aperiodic, then $\pi$ is unique and $\pi_0 P^t \to \pi$ for
\emph{every} start $\pi_0$ --- the chain forgets its origin. The rate
is spectral: ordering the eigenvalue moduli
$1 = |\lambda_1| > |\lambda_2| \ge \cdots$,
\begin{equation}
  \mathrm{TV}\bigl(\pi_0 P^t,\, \pi\bigr) \;=\; O(|\lambda_2|^t).
  \label{eq:mixing}
\end{equation}
\end{theorem}
Measured on the weather chain of the scripts: the TV decay slope is
$-0.8098$ against $\log|\lambda_2| = -0.8098$ --- the spectrum
predicting the dynamics to four decimals. The animation shows three
point-mass starts flowing to the same $\pi$.

\begin{theorem}[Ergodic theorem]
For an irreducible chain, time averages along ONE path converge:
\begin{equation}
  \frac{1}{T}\,\#\{t \le T : X_t = i\} \;\xrightarrow{\ a.s.\ }\;
  \pi_i .
  \label{eq:ergodic}
\end{equation}
\end{theorem}
This is the law of large numbers for dependent data, and it is
quietly load-bearing for the whole LLM track: it is the reason one
long text is enough to learn a chain. Measured: occupancy matches
$\pi$ to $0.002$ at $T = 2\times10^5$.

\section{Maximum likelihood on a dependent stream}

Observe $R$ paths $x^{(r)}$ of length $T$. Taking logs of
\eqref{eq:path} and collecting counts $n_{ij}$ (transitions $i \to j$
pooled over paths) and first-state counts $f_i$:
\begin{equation}
  \ell(\pi_0, P)
  \;=\; \sum_i f_i \log \pi_0(i)
  \;+\; \sum_{i}\sum_{j} n_{ij} \log P_{ij}.
  \label{eq:loglik}
\end{equation}
No term couples two rows of $P$, or $P$ with $\pi_0$: the problem
separates into $V{+}1$ independent multinoulli MLEs, each solved by
the count-and-normalize argument (Lagrange proof in
\texttt{Generative/text-to-text/bigram.tex}, Gibbs optimality in
\texttt{Theory/information-theory}):
\begin{equation}
  \boxed{\;\hat{P}_{ij} \;=\; \frac{n_{ij}}{n_i},
  \qquad
  \hat{\pi}_0(i) \;=\; \frac{f_i}{R}\;}
  \label{eq:mle}
\end{equation}
with $n_i = \sum_j n_{ij}$. The counts $(n_{ij}, f_i)$ are
\emph{sufficient}: the entire dataset enters through them.

\begin{remark}[The initial-state subtlety]
From ONE path, $\hat\pi_0$ is a point mass on the observed first state
--- degenerate, and unimprovable: a single draw from $\pi_0$ is all
the data contains about it. The LLM track silently dropped this one
factor of \eqref{eq:path} (one term in $10^6$). With many paths it
becomes an ordinary multinoulli estimate: measured, $5000$ paths
recover a non-uniform $\pi_0 = (0.5, 0.2, 0.3)$ exactly to three
decimals.
\end{remark}

\subsection{The report card, under dependence}

The samples are correlated; the bernoulli report card survives with
new proofs, every line verified by \texttt{markov\_mle.py} on a
3-state weather chain:

\paragraph{Consistent.} Ergodicity \eqref{eq:ergodic} replaces the
LLN: state $i$ is visited $n_i \approx T\pi_i \to \infty$ times, and
conditional on the visit times, row $i$'s transitions are i.i.d.\
draws from row $P_i$ (the Markov property makes each next-step draw
depend only on the current state --- the dependence structure is
exactly what makes the row samples clean). Hence
$\hat P_i \to P_i$.

\paragraph{Rate.} The row sample size is random but concentrates at
$T\pi_i$, giving
\begin{equation}
  \mathrm{Var}(\hat P_{ij})
  \;\approx\; \frac{P_{ij}\,(1 - P_{ij})}{T\,\pi_i}
  \label{eq:var}
\end{equation}
--- the multinoulli variance, weighted by occupancy: \emph{rare states
learn slowly}. Measured over $1000$ replicated paths ($T = 2000$):
the formula matches the empirical standard deviation of every entry to
within $2.7\%$. Aggregate RMS error falls with log-log slope $-0.480$
against the theoretical $-1/2$ over $T = 500 \to 128{,}000$.

\paragraph{Confidence intervals.} $\hat P_{ij} \pm 1.96
\sqrt{\hat P_{ij}(1-\hat P_{ij})/n_i}$ covers the truth in $94.7\%$ of
$1000$ replications --- the CLT behind it is a martingale CLT rather
than the i.i.d.\ one, but the practice is unchanged.

\begin{remark}[Bias, honestly]
$\hat P_{ij}$ is unbiased \emph{conditional on the row counts} $n_i$,
but $n_i$ is random and correlated with the estimates, so exact
unconditional unbiasedness fails at order $1/T$; measured:
$\max_{ij} |\mathbb{E}\hat P_{ij} - P_{ij}| = 0.0007$ at $T = 2000$.
Dependence costs a vanishing bias term, nothing more.
\end{remark}

\subsection{The fitted chain is self-consistent}

A pleasing closed loop between \eqref{eq:mle} and stationarity: for a
chain fitted to one long path, the empirical transition counts form a
flow that enters and leaves every state equally often (up to the two
path endpoints), so the letter-frequency vector is \emph{almost
exactly} stationary for $\hat P$. Measured on Shakespeare
(\texttt{Generative/text-to-text/bigram}'s $65 \times 65$ chain): the leading left
eigenvector of $\hat P$ matches the empirical character frequencies to
$1.0\times10^{-6}$ --- an eigenvector of a $65\times65$ matrix
predicting the histogram of a million characters.

\section{Experiments}

\texttt{markov\_chains\_demo.py} (probability): Chapman--Kolmogorov,
ergodic occupancy, the mixing animation (three starts flowing to one
$\pi$, TV decay on the $|\lambda_2|^t$ line), the $k$-th order
reduction on Shakespeare trigrams, and the stationary-vs-histogram
check. \texttt{markov\_mle.py} (statistics): the $\sqrt{T}$ law on
dependent data, the occupancy-weighted variance formula
\eqref{eq:var}, CI coverage, $\hat\pi_0$ from many paths, a 14-day
forecast sampled from the \emph{learned} weather chain, and the
animation --- one path streaming in, $|\hat P - P|$ fading, the error
descending the $-1/2$ line, occupancy bars converging to $\pi$.

One sentence to keep: \emph{dependence does not break maximum
likelihood --- it only makes the sample sizes random, with the ergodic
theorem standing where the law of large numbers used to.}

\end{document}
